\chapter{Grundlagen}
\label{ch:02_grundlagen}

In diesem Kapitel werden die theoretischen Fundamente der untersuchten Rendering-Pipelines definiert.
Dazu gehört die Abgrenzung der verschiedenen Render-Pfade, die Spezifikation der physischen Hardware-Flaschenhälse sowie die Einordnung der Konzepte im Kontext der Stilisierung.


\section{Rendering-Paradigmen und algorithmische Komplexität}
\label{sec:0201_rendering_paradigms_and_complexity}

Dieser Abschnitt widmet sich der formalen Definition und Gegenüberstellung von klassischen und entkoppelten Render-Pfaden unter Berücksichtigung ihrer inhärenten Skalierbarkeit.

\subsection{Forward Rendering}
\label{subsec:020101_forward_rendering_core}
Beschreibung der klassischen Render-Pipeline mit der direkten Berechnung von Geometrie und Schattierung in einem Durchlauf. Erläuterung des theoretischen Aufwandsverhaltens bei steigender Objekt- und Lichtanzahl.

\subsection{Deferred Rendering}
\label{subsec:020102_deferred_rendering_core}
Erklärung der entkoppelten Pipeline und der Zwischenspeicherung von Geometriedaten in separaten Puffern. Aufzeigen der Effizienzgewinne bei der Schattierung durch den Einsatz geometrischer Hilfskörper.


\section{Hardware-Limitierungen und Performance-Metriken}
\label{sec:0202_hardware_limitations_and_metrics}

Dieser Abschnitt spezifiziert die physischen Ausführungseinheiten moderner Grafikhardware, welche die kritischen Engpässe bei der praktischen Ausführung der Render-Pfade bilden.

\subsection{Fragment-Overdraw und Fillrate}
\label{subsec:020201_fragment_overdraw_and_fillrate}
Thematisierung mehrfach gezeichneter Pixel auf identischen Koordinaten. Betrachtung der resultierenden Last für die Füllrate und softwareseitiger Optimierungen wie dem Z-Prepass.

\subsection{Speicherbandbreite (Memory Bandwidth)}
\label{subsec:020202_memory_bandwidth_limitations}
Analyse der Durchsatzgrenzen beim Zugriff auf den Videospeicher. Erläuterung des Lese- und Schreibaufwands durch breite Pufferstrukturen und additive Mischvorgänge.

\subsection{Register Pressure}
\label{subsec:020203_register_pressure_limitations}
Behandlung der limitierten internen Speichereinheiten pro Rechenkern. Beschreibung der Auswirkungen auf die Thread-Parallelisierung bei der Nutzung großer, kombinierter Shader-Programme.

\subsection{Dynamic Branching und Warp Divergence}
\label{subsec:020204_dynamic_branching_and_warp_divergence}
Erläuterung der synchronen Funktionsweise von Thread-Gruppen auf Hardware-Ebene. Beschreibung der sequenziellen Serialisierung bei ungleichen Verzweigungsbedingungen auf Pixelebene.


\section{Grundlagen der Non-Photorealistic Rendering (NPR) Pipeline}
\label{sec:0203_non_photorealistic_rendering_fundamentals}

Dieser Abschnitt kategorisiert die raumbezogenen Konzepte der NPR-Generierung zur selektiven, isolierten Stilisierung von Szenenobjekten.

\subsection{Object-Space- vs. Screen-Space-Stilisierung}
\label{subsec:020301_object_space_vs_screen_space_stylization}
Abwägung von Stilisierungsverfahren direkt auf der dreidimensionalen Szenengeometrie im Vergleich zu nachgelagerten, globalen Filtern auf dem fertigen zweidimensionalen Bild. Gegenüberstellung von präziser Objekttrennung und Laufzeit-Konstanz.